(1) 极坐标轴对称振动满足 \(\rho_s u_{tt}=S[u_{rr}+(1/r)u_r]\)，令 \(c=\sqrt{S/\rho_s}\)。固定边缘条件是 \(u(R,t)=0\)，中心 \(r=0\) 处位移与导数必须有限；不允许把在原点发散的 Neumann 函数 \(Y_0\) 作为物理振型。
(2) 取分离变量 \(u(r,t)=f(r)e^{i\omega t}\)，径向方程变成 \(f''+f'/r+k^2f=0\)，其中心正则解为 \(f=A J_0(kr)\)。边界 \(J_0(kR)=0\) 给 \(k_n=j_{0,n}/R\)，因此 \(\omega_n=c j_{0,n}/R\)。正交关系为 \(\int_0^RrJ_0(j_{0,n}r/R)J_0(j_{0,m}r/R)\,dr=(R^2/2)J_1(j_{0,n})^2\delta_{nm}\)。
(3) 利用前述正交投影，\(A_n=[2u_0/J_1(j_{0,n})^2]\int_0^1 s(1-s^2)J_0(j_{0,n}s)\,ds\)。在 \(J_0(j_{0,n})=0\) 时积分等于 \(4J_1(j_{0,n})/j_{0,n}^3\)，因此 \(A_n=8u_0/[j_{0,n}^3J_1(j_{0,n})]\)。由于初速度为零，完整解为 \(u(r,t)=\sum_{n\ge1}A_nJ_0(j_{0,n}r/R)\cos(\omega_nt)\)，满足中心正则与固定边缘。
(4) 初始动能为零，初始位移梯度为 \(u_r=-2u_0r/R^2\)。应变能积分 \(E=(S/2)\int_0^{2\pi}\int_0^R (u_r)^2r\,dr\,d\phi\) 得 \(E=\pi S u_0^2\)。不同模的能量交叉项由正交性消去，每个无耗散模分别守恒总动能加势能，所以模态级数对任意 \(t\) 的总能仍为 \(\pi S u_0^2\)。
